\section{AMI Labs：组织形式也是技术路线的一部分}

当访谈进入 AMI Labs，问题从“世界模型是什么”变成“什么组织能做世界模型”。谢赛宁认为，在学校里资源不足，在大厂里产品周期和榜单军备竞赛会压缩探索空间；传统研究院有自由但难以承担大规模训练和真实场景合作；完全封闭的大模型公司又可能切断学术连接和开放讨论。因此，AMI Labs 试图找到一种中间形态。

\begin{figure}[H]
\centering
\includegraphics[width=0.96\textwidth]{ami-labs-thesis.png}
\caption{AMI Labs 的组织命题：在世界模型主线、研究友好组织、全球伙伴网络、商业闭环和学术连接之间找平衡。}
\end{figure}

\begin{knowledgebox}{读图：组织设计如何影响技术可能性}
图中五个节点共同构成 AMI 的命题。世界模型需要长期 research；长期 research 需要资源和组织氧气；真实世界数据需要全球伙伴网络；严肃创业需要商业闭环；研究路线又需要和学术界保持连接。任何一个节点缺失，world model 都可能退化成论文项目、封闭产品或空泛愿景。
\end{knowledgebox}

\subsection{不是纯研究院，也不是封闭产品公司}

前面的图已经给出 AMI 的五个节点；本节先看最核心的组织张力。世界模型需要研究自由，但创业公司又必须处理资源、客户、资本和交付压力。

谢赛宁说 AMI 既不是 non-profit，也不是纯粹 research lab；它需要 business model。但它也不想成为只围绕产品 deadline 和 benchmark 打转的封闭公司。这个张力很难，因为商业公司必须对资源和结果负责，而世界模型又是高度探索性的方向。

这也是为什么他强调“researcher friendly organization”。如果组织没有足够氧气，研究者即使知道某个问题重要，也会被安排去做短期产品链条里可交付的环节，比如 video captioning、榜单优化或发布周期支持。

\begin{importantbox}{组织形式会改变研究空间}
同一个研究者、同一个 idea，在学校、大厂、封闭创业公司和研究友好创业公司里，能做的事情完全不同。世界模型这样的大问题，既需要资源，又需要探索自由，还需要真实场景合作。
\end{importantbox}

\subsection{Yann LeCun 的作用：路线、人格与信任}

组织张力讲完后，接下来要看谁来承载这种张力。访谈里 Yann LeCun 不只是共同创始人的名字，而是路线信念、科学家 integrity 和团队心理稳定器的结合。

访谈中，谢赛宁多次谈到 Yann LeCun。他强调 LeCun 并不反对 LLM，而是反对“LLM 自然通向人类级智能”的叙事。他敬佩 LeCun 的地方，一是长期坚持 world model / JEPA / autonomous intelligence 的路线，二是科学家的 integrity，三是个人气质：热爱生活、有广阔世界、能让身边的人感到前方有路。

这段内容对技术组织有启发：founder 或科学领袖不仅提供方向，还提供组织心理。研究问题越不确定，团队越需要有人提供长期信念、及时校正和允许质疑的讨论空间。

\subsection{本章小结}

AMI Labs 的故事说明，技术路线和组织路线不可分。世界模型不是只靠一个模型结构推进的，它需要资源、伙伴、研究文化、商业闭环和长期叙事共同支撑。

